\subsection{引理 A\texorpdfstring{$''$}{} 的初等修订}\label{sec:lemAdp}
\begin{theorem}[引理 A$''$, 带正余量]\label{thm:lemAdp}
设 $R\ge1500$, $0<u<1/2$, $w=u\sqrt R\ge2$. 则
\begin{equation}
 G(R,u)>\bar D(u)+\frac{\ell}{Ru^3}>\bar D(u).
\end{equation}
\end{theorem}
\begin{proof}
令 $c=\varepsilon\ell/u$, $\alpha=\varepsilon c=\ell/(Ru)$,
$t=a(u)$, $\theta=\theta_2$, $v=u/\ell=-t\cot t>0$.
由前节的模式识别, $0<\theta_1<\pi/2<\theta_2<\pi$, 所有下述三角除法合法.
还有 $0<c\le1/4-\varepsilon$, $\varepsilon<1/38$, $\alpha<1/152$.
奇方程与 $\cot z<1/z$ 给出
\[
 \delta_2\le\varepsilon\cot z_2<\frac{\varepsilon}{c\theta_2}
 \le\frac{2\varepsilon}{\pi c},\qquad
 z_2<\frac\pi8-\varepsilon(\pi/2-2/\pi)<\frac\pi8.
\]
函数 $s\mapsto-s\cot s$ 在 $(\pi/2,\pi)$ 上严格递增, 而
$-\theta\cot\theta=\varepsilon\theta\cot(c\theta)<\varepsilon/c=v=-t\cot t$,
所以 $\theta<t$. 定义
\[
 d_1=\pi^2/4-\theta_1^2>0,\qquad d_2=t^2-\theta^2>0,\qquad
 G-\bar D=(d_1-d_2)/u^2.
\]

\emph{第一差量下界.} $z_1<\pi/8$ 且 $\tan(\pi/8)=\sqrt2-1<1/2$,
故 $0<\delta_1<1/76$ 及 $\theta_1>7/5$.
用 $\tan z\ge z$, $\arctan x\ge x-x^3/3$, $\pi^2<10$ 得
\[
 \delta_1\ge\arctan(\alpha\theta_1)
 \ge\alpha\theta_1(1-\pi^2\alpha^2/12)
 >\frac{99}{100}\alpha\theta_1.
\]
因此
\begin{equation}
 d_1=(\pi/2+\theta_1)\delta_1
 >\frac{29}{10}\frac75\frac{99}{100}\alpha
 =\frac{20097}{5000}\alpha>4\alpha.\label{eq:r8-def1}
\end{equation}

\emph{余切余项与一个连续三角界.} 令 $r(z)=1/z-\cot z$.
由 $\sin z-z\cos z=\int_0^z s\sin s\,ds$,
$\sin s\le s$ 与 $\sin z\ge z-z^3/6$, 对 $0<z\le\pi/8$ 有
\begin{equation}
 0<r(z)\le\frac{z}{3(1-z^2/6)}
 <\frac{64}{187}z<\frac7{20}z.\label{eq:r8-cot-elementary}
\end{equation}
这里仅用 $3<\pi<22/7$, 因而 $\pi^2<10$.
另对 $\pi/2<t<\pi$,
\begin{equation}
 B(t):=\frac{2t^4}{t^2+v^2+v}
 =\frac{2t^3\sin^2t}{t-\sin t\cos t}<2(t\sin t)^2\le8.\label{eq:r8-B8}
\end{equation}
最后一步可直接证明: $t\le2$ 时 $t\sin t\le2$.
$t\ge2$ 时 $q(t)=t\sin t$ 满足 $q''=2\cos t-t\sin t<0$.
交错 Taylor 余项给出
\[
 \sin2\le\frac{2578}{2835}<\frac{91}{100},\qquad
 \cos2\le-\frac{131}{315}<-\frac25.
\]
故凹函数的切线估计与 $t-2<8/7$ 给出
\[
 q(t)\le q(2)+q'(2)(t-2)
 <\frac{91}{50}+\frac{11}{100}\frac87
 =\frac{681}{350}<2.
\]

\emph{第二差量上界.} 写 $\psi=t-\theta\in(0,\pi/2)$,
$A=\tan(t-\pi/2)=v/t$, $B=\tan(\theta-\pi/2)=\varepsilon\cot(c\theta)$.
恒等式
\[
 A-B=-\frac{v\psi}{t\theta}+\varepsilon r(c\theta),\qquad
 \tan\psi=\frac{A-B}{1+AB}
\]
结合 $\tan\psi\ge\psi$ 给出
\[
 \psi\le\frac{\varepsilon r(c\theta)}{1+AB+v/(t\theta)}.
\]
取 $C=7/20$, 由 \eqref{eq:r8-cot-elementary} 有
\[
 \varepsilon r(c\theta)<C\alpha\theta,\qquad
 AB>\frac{v^2}{t\theta}-C\varepsilon^2\frac\theta t.
\]
令 $D_0=1+v(v+1)/(t\theta)\ge1$, $d=C\varepsilon^2\theta/t<1/1000$.
固定 $t,v$ 时, $\theta(t+\theta)/D_0$ 的正分子随 $\theta$ 增加,
正分母随 $\theta$ 减少, 因而整个比值增加. 由 $\theta<t$ 和 \eqref{eq:r8-B8},
\begin{equation}
 d_2=(t+\theta)\psi
 < C\alpha\frac{\theta(t+\theta)}{D_0-d}
 <\frac7{20}\frac{1000}{999}\,8\alpha
 =\frac{2800}{999}\alpha<3\alpha.\label{eq:r8-def2}
\end{equation}
所有分母均有上述正下界. 由 \eqref{eq:r8-def1} 和 \eqref{eq:r8-def2},
$G-\bar D>\alpha/u^2=\ell/(Ru^3)$, 得证.
\end{proof}

本修订主链仅需上述初等粗界; 特别地 $d_2/d_1<3/4<0.8256$.
旧主链所用更细的 $C_z<0.337$ 和 $B(t)\le9$ 另经第八轮精确证书修复,
作为可复用的标量工具保留, 不再承担 T1 的连续区域覆盖.
